\section{A dihedral quotient gives the type restriction}\label{a-dihedral-quotient-gives-the-type-restriction}

\textbf{Lemma 5.} If a finite group \(G\) of order \(2^n p\) is generated by
involutions, then \(G/O_2(G)\cong D_{2p}\), where \(O_2(G)\) is its largest
normal two-subgroup.

\emph{Proof.} Put \(\overline G=G/O_2(G)\). It still has an odd part of order
\(p\), is generated by involutions or identities, and has no nontrivial
normal two-subgroup. Burnside's theorem makes it solvable. A minimal
normal subgroup is therefore elementary abelian, and must have order
\(p\). Write it as \(P\).

Its centralizer \(C=C_{\overline G}(P)\) is normal. If \(S\) is a Sylow
two-subgroup of \(C\), then \(C=P\times S\): \(P\) is central in \(C\) and the
orders multiply to \(|C|\). Thus \(S\) is the unique Sylow two-subgroup of
\(C\), hence characteristic there and normal in \(\overline G\). It is
trivial. Therefore \(C=P\) and
\(\overline G/P\hookrightarrow\operatorname{Aut}(P)\) is cyclic.
Because it is generated by elements of order at most two, this quotient
has order at most two. It cannot be trivial: the nontrivial group of
odd order \(P\) cannot be generated by involutions. Its order is two and
its action on \(P\) is the unique nontrivial involution of
\(\operatorname{Aut}(C_p)\), namely inversion. A generating involution
outside \(P\) supplies a complement. Hence \(\overline G=D_{2p}\).
\(\square\)

Suppose neither adjacent product has order divisible by \(p\). Their
images in \(D_{2p}\) have order at most two, and the endpoint product also
has order at most two. Three involutions or identities with all pair
products of order at most two commute pairwise. They would generate an
abelian two-group, not \(D_{2p}\). This contradiction proves Theorem 1.
The allowed type forms follow from Lagrange's theorem. The exponents
\(s,t\) in this type restriction are not silently restricted to those of
the narrower construction in Theorem 2.
